\section{Descripción del Problema}

ECOBICI es un sistema de movilidad urbana basado en el préstamo y renta de bicicletas dentro de la Ciudad de México. El servicio permite a los usuarios adquirir diferentes tipos de membresías para acceder a las bicicletas disponibles en diversas estaciones distribuidas estratégicamente en la ciudad.

Debido al volumen de información generado por las operaciones diarias, resulta necesario contar con una base de datos capaz de administrar usuarios, membresías, métodos de pago, estaciones, bicicletas, viajes, incidentes y personal de la organización. Además, el sistema debe permitir mantener un historial de viajes, registrar accidentes e incidentes, controlar el inventario de bicicletas y estaciones, gestionar la información de los empleados y generar reportes estadísticos que apoyen la toma de decisiones.

Por lo anterior, se requiere diseñar e implementar una base de datos relacional que garantice la integridad, consistencia y disponibilidad de la información, permitiendo automatizar procesos operativos y administrativos mediante el uso de restricciones, procedimientos almacenados, funciones, vistas y disparadores.
